\documentclass[kdmile,a4paper]{kdmile}

\usepackage{graphicx,url}
\usepackage[T1]{fontenc}
\usepackage{booktabs}
\usepackage{multirow}

% Teoremas e definições
\newtheorem{theorem}{Theorem}[section]
\newtheorem{conjecture}[theorem]{Conjecture}
\newtheorem{corollary}[theorem]{Corollary}
\newtheorem{proposition}[theorem]{Proposition}
\newtheorem{lemma}[theorem]{Lemma}
\newdef{definition}[theorem]{Definition}
\newdef{remark}[theorem]{Remark}

% ============================================================
% CABEÇALHO: ajuste nomes conforme número de autores
% 2-3 autores: iniciais; >3 autores: "A. Sobrenome et al."
% ============================================================
\markboth{E. Barreto and K. S. Lima}
         {Price Disparity between Healthy and Processed Foods as a Predictor of Obesity Rates in the USA}

% ============================================================
% TÍTULO
% ============================================================
\title{Price Disparity between Healthy and Processed Foods as a\\
       Predictor of Obesity Rates in the USA:\\
       An Analysis with USDA and BRFSS Public Data}

% ============================================================
% AUTORES E AFILIAÇÕES
% ============================================================
\author{Evandro Barreto\inst{1} \and Kaiky de Souza Lima\inst{1}}

\institute{
  % Substitua pela instituição real de vocês
  Universidade Tecnológica Federal do Paran\'{a}, Brazil \\
  \email{\{evandrobarreto, kaiky\}@alunos.utfpr.edu.br}
}

% ============================================================
% ABSTRACT (obrigatório em inglês, 100–300 palavras, 1 parágrafo)
% ============================================================
\begin{abstract}
Obesity is one of the most pressing global public health challenges, with socioeconomic
factors playing a decisive role in its prevalence. This study investigates whether the
price disparity between healthy and ultra-processed foods is a statistically significant
predictor of obesity rates across American states. We integrate two public datasets: the
Behavioral Risk Factor Surveillance System (BRFSS), provided by the CDC, which contains
time-series indicators of obesity, fruit and vegetable consumption, and physical activity
among American adults; and the Food-at-Home Monthly Area Prices dataset, provided by the
USDA, which monitors regional monthly prices of household food items disaggregated by
geographic area. A two-stage predictive analytical pipeline is proposed: first, a
classification model to assess whether variations in healthy food prices significantly
influence obesity rates; second, conditional on that relationship, a regression model to
quantify the magnitude of this impact. Additionally, income level is examined as a
potential moderating factor in the relationship between food prices and obesity risk.
We hypothesize that sustained increases in fruit and vegetable prices are positively
associated with higher obesity rates, and that lower-income populations are
disproportionately affected. This work is aligned with the topics of Classification,
Supervised Learning, Machine Learning Applications, Data Mining Applications, and
Time-Series Analysis.
\end{abstract}

% ============================================================
% CCS CONCEPTS — gere em: https://dl.acm.org/ccs/ccs.cfm
% ============================================================
\begin{CCSXML}
<ccs2012>
  <concept>
    <concept_id>10010147.10010257</concept_id>
    <concept_desc>Computing methodologies~Machine learning</concept_desc>
    <concept_significance>500</concept_significance>
  </concept>
  <concept>
    <concept_id>10010147.10010257.10010293</concept_id>
    <concept_desc>Computing methodologies~Supervised learning</concept_desc>
    <concept_significance>300</concept_significance>
  </concept>
</ccs2012>
\end{CCSXML}

\ccsdesc[500]{Computing methodologies~Machine learning}
\ccsdesc[300]{Computing methodologies~Supervised learning}

\keywords{obesity, food prices, machine learning, classification, supervised learning,
          BRFSS, USDA, data mining}

% ============================================================
% FINANCIAMENTO (se houver)
% ============================================================
% \begin{bottomstuff}
% This work was partially funded by CNPq grant number XYZ.
% \end{bottomstuff}

\begin{document}
\begin{bottomstuff}
\end{bottomstuff}
\maketitle

% ============================================================
% 1. INTRODUÇÃO
% ============================================================
\section{Introduction}

A obesidade constitui um dos maiores desafios contemporâneos de saúde pública, sendo
substancialmente agravada por variáveis macroeconômicas.

Embora a relação entre custo alimentar e obesidade seja intuitivamente sugerida na
literatura, a principal lacuna metodológica reside na integração sistemática de bases de
dados epidemiológicas e econômicas díspares para atestar esse fenômeno de forma
quantitativa e replicável. Este artigo aborda essa lacuna por meio de um pipeline analítico
preditivo que combina dados públicos do USDA e do CDC/BRFSS.

As seguintes perguntas de pesquisa orientam o estudo:

\begin{itemize}
  \item \textbf{P1:} Existe relação mensurável e estatisticamente significativa entre o
  custo de alimentos saudáveis (preço de frutas e vegetais) e a faixa de renda da
  população com os índices de obesidade nos estados norte-americanos?
\end{itemize}

As hipóteses testadas são:
\begin{itemize}
  \item \textbf{H1:} A elevação sustentada nos preços médios de frutas e vegetais está
  positivamente associada ao aumento das taxas de obesidade.
  \item \textbf{H2:} Categorias específicas de grupos de alimentos apresentam variações de
  preço com impacto diferenciado nos índices de obesidade.
  \item \textbf{H3:} Populações com menor faixa de renda apresentam maiores índices de
  obesidade, sendo a renda um fator moderador na relação entre custo de alimentos
  saudáveis e risco de obesidade.
\end{itemize}

% ============================================================
% 2. TRABALHOS RELACIONADOS
% ============================================================
\section{Related Work}
\label{sec:related}

% TODO: expandir com referências reais
Em desenvolvimento \cite{TODO} 


% ============================================================
% 3. DADOS
% ============================================================
\section{Data}
\label{sec:data}

\subsection{Dataset 1: Nutrition, Physical Activity and Obesity (CDC/BRFSS)}

O Behavioral Risk Factor Surveillance System (BRFSS) é um sistema de vigilância em saúde
conduzido pelo CDC (Centers for Disease Control and Prevention). O dataset utilizado
contém séries temporais e recortes geográficos com indicadores de obesidade, consumo de
frutas e vegetais e atividade física de adultos americanos.

As principais variáveis de interesse incluem: taxa de obesidade por estado e ano,
estratificada por faixa de renda, nível educacional, sexo e etnia.

\subsection{Dataset 2: Food-at-Home Monthly Area Prices (USDA)}

O dataset do USDA monitora mensalmente os preços regionais de gêneros alimentícios
domiciliares, desagregado por área geográfica. Permite análise temporal das oscilações de
custo e cruzamento geoespacial com os dados de saúde do BRFSS.

As principais variáveis de interesse incluem: preço médio mensal por categoria de alimento
(frutas, vegetais, proteínas, processados) por região geográfica.

\subsection{Integração dos Dados}

A integração entre os dois datasets é realizada por meio de chave composta
\texttt{(estado, ano)}, permitindo o alinhamento temporal e geoespacial entre as
variáveis de preço e os indicadores de saúde. A Tabela~\ref{tab:datasets} sumariza as
características principais dos datasets.

\begin{table}[h!]
\caption{Características dos datasets utilizados}
\label{tab:datasets}
\centering
\begin{tabular}{lll}
\toprule
\textbf{Atributo}  & \textbf{BRFSS (CDC)}          & \textbf{FAHMAP (USDA)} \\
\midrule
Granularidade      & Estado / ano                  & Região / mês           \\
Período            & 2012--2018                    & 2012--2018             \\
Nº de registros    & 110.880                       & 1.020.600              \\
Variável-alvo      & Taxa de obesidade (\%)        & --                     \\
\bottomrule
\end{tabular}
\end{table}

% ============================================================
% 4. METODOLOGIA
% ============================================================
\section{Methodology}
\label{sec:method}

O pipeline analítico é composto por três etapas: pré-processamento e integração dos
dados, análise exploratória e correlacional, e — como trabalho futuro — modelagem
preditiva supervisionada.

\subsection{Pré-processamento e Integração}

Os dados do BRFSS foram filtrados para registros com indicador de obesidade válido,
estratificados por faixa de renda e agregados por estado e ano. Os dados do USDA foram
agregados de granularidade mensal para médias anuais por região geográfica. A integração
entre os dois datasets foi realizada por meio de chave composta \texttt{(região, ano)},
resultando em um dataset analítico unificado com variáveis de preço e indicadores de
saúde alinhados temporalmente. Todo o pipeline foi implementado em Python utilizando
\texttt{pandas} e \texttt{DuckDB} para consultas SQL sobre os dados brutos.

\subsection{Análise Exploratória e Correlacional}

Foram construídas duas variáveis derivadas principais: \textit{Price\_Healthy}, o preço
médio de alimentos saudáveis (frutas e vegetais), e \textit{Price\_Ratio}, a razão entre
o preço de alimentos saudáveis e ultraprocessados. A relação entre essas variáveis e as
taxas de obesidade foi investigada por meio do coeficiente de correlação de Pearson,
calculado separadamente para cada região geográfica (Nordeste, Centro-Oeste, Sul e
Oeste). A influência da faixa de renda foi avaliada por meio de correlação entre uma
codificação ordinal das faixas de renda e as taxas de obesidade, além de análise
temporal estratificada por grupo de renda.

\subsection{Plano de Modelagem Preditiva (Trabalho Futuro)}

Em desenvolvimento


% ============================================================
% 5. PLANO DE AVALIAÇÃO EXPERIMENTAL
% ============================================================
\section{Experimental Evaluation Plan}
\label{sec:experiments}

% TODO: substituir pelo texto de resultados quando disponíveis

\subsection{Configuração Experimental}

Em desenvolvimento

\subsection{Métricas de Avaliação}

Em desenvolvimento

\subsection{Hipóteses e Experimentos Planejados}

Em desenvolvimento

\subsection{Resultados Preliminares}

% TODO: inserir resultados da análise exploratória aqui
Em desenvolvimento


% ============================================================
% 6. RESULTADOS E DISCUSSÃO
% ============================================================
\section{Results and Discussion}
\label{sec:results}

Em desenvolvimento.

% ============================================================
% 7. LIMITAÇÕES E AMEAÇAS À VALIDADE
% ============================================================
\section{Limitations and Threats to Validity}
\label{sec:limitations}

Em desenvolvimento.

% ============================================================
% 8. CONCLUSÃO
% ============================================================
\section{Conclusion}
\label{sec:conclusion}

% TODO: preencher após os experimentos
Em desenvolvimento

% ============================================================
% AGRADECIMENTOS (opcional)
% ============================================================
% \begin{ack}
% Os autores agradecem...
% \end{ack}

\begin{received}
\end{received}

% ============================================================
% REFERÊNCIAS
% ============================================================
\bibliographystyle{kdmile}
\bibliography{kdmileb}

% Referências manuais temporárias (remover ao usar .bib):
% \begin{thebibliography}{9}
%
% \bibitem{TODO}
% Autor, A. A. e Autor, B. B.
% Título do artigo.
% In: \textit{Proceedings of ...}, pp. XX--XX, Ano.
%
% \end{thebibliography}

\end{document}